\chapter{Analyse et Spécification des Besoins}

\section{Introduction}
L'analyse des besoins constitue la fondation de notre projet. Elle permet de traduire les attentes des utilisateurs en spécifications techniques rigoureuses.

\section{Analyse Fonctionnelle Détaillée des Interfaces}
eduGenius est structuré autour de trois espaces de travail distincts, chacun répondant à des besoins métiers spécifiques.

\subsection{L'Espace Administrateur : Le Cerveau Académique}
L'interface administrateur d'eduGenius a été conçue pour couvrir l'intégralité du cycle de gestion d'une université moderne.
\begin{itemize}
    \item \textbf{Gestion des Périodes Académiques} : Définition des années universitaires et des semestres avec leurs dates de début et de fin respectives.
    \item \textbf{Module de Configuration LMD} : Gestion des programmes d'études (Licence, Master, Ingénierie) et des matières avec paramétrage fin des coefficients et crédits.
    \item \textbf{Centre de Visioconférence (Live Monitoring)} : Surveillance en temps réel des sessions de cours actives, avec capacité de terminer une session à distance pour forcer le calcul de présence.
    \item \textbf{Reporting d'Assiduité} : Génération de rapports statistiques sur le taux de présence moyen, avec filtres par période et export au format CSV pour les archives de Polytech-Intl.
    \item \textbf{Module de Génération de PV} : Un outil critique permettant de générer les \textbf{Procès-Verbaux (PV)} officiels des examens, incluant les moyennes calculées, les mentions automatiques et les statistiques de réussite par session (Principale, Rattrapage).
\end{itemize}

\subsection{L'Espace Enseignant : L'Assistant Pédagogique IA}
\begin{itemize}
    \item \textbf{Gestion des Travaux et Rendus} : Organisation des supports par chapitres et création d'exercices avec dates limites. Suivi des fichiers déposés par les étudiants.
    \item \textbf{Système de Correction Avancé} : Interface de notation par question avec ajout de feedbacks textuels spécifiques pour chaque erreur détectée.
    \item \textbf{Studio de Quiz IA} : Création de quiz manuels ou assistés par l'IA (Gemini), avec gestion de la difficulté et du chronométrage.
    \item \textbf{Lobby de Session Live} : Interface de préparation avant le lancement d'une salle Daily.co, affichant les détails de la matière et le lieu virtuel.
\end{itemize}

\subsection{L'Espace Étudiant : Gamification et Apprentissage}
L'étudiant bénéficie d'une interface motivante intégrant des éléments de \textbf{Gamification}.
\begin{itemize}
    \item \textbf{Système d'Expérience (XP)} : Chaque quiz complété et chaque cours consulté rapporte des points d'expérience, permettant de suivre son engagement de manière ludique.
    \item \textbf{Centre IA Personnel} : Accès direct aux générateurs de flashcards, résumés et quiz d'entraînement basés sur l'historique de performance.
    \item \textbf{Dashboard de Performance} : Visualisation de la Moyenne Générale (GPA) estimée et des moyennes par matière avec indicateurs de tendance.
    \item \textbf{Agenda Dynamique} : Vue calendrier hebdomadaire intégrant les cours en présentiel (numéros de salles) et les cours en ligne (liens de réunion).
\end{itemize}

\section{Spécification Détaillée des Cas d'Utilisation}
...

\subsection{Cas d'Utilisation : Gestion des Utilisateurs}
\begin{table}[H]
\centering
\begin{tabular}{|l|p{10cm}|}
\hline
\textbf{Nom du Cas} & \textbf{Gérer les Comptes Utilisateurs} \\ \hline
\textbf{Acteur} & Administrateur \\ \hline
\textbf{Pré-conditions} & L'administrateur est authentifié. \\ \hline
\textbf{Scénario Nominal} & 
1. L'admin accède à la liste des membres. \newline 
2. Il choisit de créer, modifier ou supprimer un compte. \newline
3. Le système valide les données (email unique, format du MDP). \newline
4. La base de données MongoDB est mise à jour. \\ \hline
\textbf{Post-conditions} & L'utilisateur reçoit ses identifiants par email (optionnel). \\ \hline
\end{tabular}
\end{table}

\subsection{Cas d'Utilisation : Publication d'Annonces}
\begin{table}[H]
\centering
\begin{tabular}{|l|p{10cm}|}
\hline
\textbf{Nom du Cas} & \textbf{Publier une Annonce de Classe} \\ \hline
\textbf{Acteur} & Enseignant \\ \hline
\textbf{Pré-conditions} & L'enseignant est assigné à la classe concernée. \\ \hline
\textbf{Scénario Nominal} & 
1. L'enseignant accède au module Annonces. \newline 
2. Il rédige le contenu et choisit s'il doit être épinglé. \newline
3. Le système enregistre l'annonce et émet un événement Socket.io. \newline
4. L'annonce apparaît sur le dashboard des étudiants en temps réel. \\ \hline
\textbf{Post-conditions} & L'annonce est persistée et visible par tous les membres de la classe. \\ \hline
\end{tabular}
\end{table}

\subsection{Cas d'Utilisation : Génération de Résumé IA}
\begin{table}[H]
\centering
\begin{tabular}{|l|p{10cm}|}
\hline
\textbf{Nom du Cas} & \textbf{Générer un Résumé Automatique} \\ \hline
\textbf{Acteur} & Enseignant \\ \hline
\textbf{Pré-conditions} & Le document PDF est uploadé sur S3. \\ \hline
\textbf{Scénario Nominal} & 
1. L'enseignant sélectionne un cours. \newline 
2. Il clique sur "Générer Résumé". \newline
3. Le système extrait le texte du PDF. \newline
4. L'IA Gemini traite le texte et renvoie le résumé. \newline
5. Le résumé est affiché et enregistré. \\ \hline
\textbf{Post-conditions} & Le résumé est accessible par tous les étudiants de la classe. \\ \hline
\end{tabular}
\end{table}

\subsection{Cas d'Utilisation : Session Live Vidéo}
\begin{table}[H]
\centering
\begin{tabular}{|l|p{10cm}|}
\hline
\textbf{Nom du Cas} & \textbf{Démarrer un cours en visioconférence} \\ \hline
\textbf{Acteur} & Enseignant \\ \hline
\textbf{Pré-conditions} & L'enseignant est connecté à son dashboard. \\ \hline
\textbf{Scénario Nominal} & 
1. L'enseignant choisit sa classe. \newline 
2. Il clique sur "Démarrer Session Live". \newline
3. Le backend demande un token à Daily.co. \newline
4. La salle virtuelle s'ouvre dans l'interface. \newline
5. Une notification est envoyée aux étudiants. \\ \hline
\textbf{Post-conditions} & La présence des étudiants est enregistrée automatiquement. \\ \hline
\end{tabular}
\end{table}

\section{Modélisation des Diagrammes de Séquences}

\subsection{Diagramme de Séquence : Authentification JWT}
Le processus d'authentification est vital pour la sécurité d'eduGenius. Il se déroule comme suit :
\begin{enumerate}
    \item L'utilisateur envoie ses identifiants (Email/Password) au serveur Express.
    \item Le serveur vérifie le mot de passe via \texttt{bcrypt.compare()}.
    \item Si valide, le serveur génère un \textbf{JSON Web Token} signé avec une clé secrète.
    \item Le token est renvoyé au client qui le stocke en \textit{LocalStorage}.
    \item Pour chaque requête ultérieure, le client inclut le token dans le header \texttt{Authorization}.
\end{enumerate}
\fbox{Placeholder: Diagramme de Séquence - Authentification JWT}

\subsection{Diagramme de Séquence : Flux de Génération IA}
Ce diagramme illustre la complexité de l'interaction avec les services tiers (AWS et Google) :
\begin{enumerate}
    \item L'utilisateur initie la demande depuis le Frontend.
    \item Le Backend récupère l'URL du fichier sur AWS S3.
    \item Le Backend extrait le texte brut via \texttt{pdf-parse}.
    \item Le texte est envoyé à l'orchestrateur \texttt{AiService} qui communique avec l'API Gemini.
    \item Le résultat JSON est parsé, validé, puis enregistré dans MongoDB.
    \item Le résultat final est renvoyé au Frontend pour affichage.
\end{enumerate}
\fbox{Placeholder: Diagramme de Séquence - Génération IA}

\section{Besoins Non-Fonctionnels}
Les besoins non-fonctionnels définissent les contraintes de qualité et les exigences de performance du système eduGenius.

\subsection{Performance et Scalabilité}
L'API doit répondre en moins de 200ms pour les opérations standards. L'utilisation de \textbf{Mongoose virtuals} et d'indexations sur MongoDB garantit une recherche rapide même avec des milliers d'utilisateurs. Pour le frontend, l'utilisation de \textbf{SWR} pour le caching des données permet une navigation instantanée.

\subsection{Disponibilité et Résilience (Le rôle d'Ollama)}
La plateforme doit garantir un service minimum même en cas d'indisponibilité des services Cloud externes. L'intégration de \textbf{Ollama} avec le modèle \textbf{Gemma 3:4b} permet de traiter les demandes urgentes de résumé localement sur les serveurs de BES2I, assurant ainsi une résilience de 99.9\% face aux pannes internet.

\subsection{Sécurité et Confidentialité (RGPD)}
Le respect de la vie privée des étudiants est primordial, conformément aux standards du RGPD. 
\begin{itemize}
    \item \textbf{Chiffrement} : Toutes les communications transitent par \textbf{HTTPS}. Les mots de passe sont hachés via \textbf{bcrypt}.
    \item \textbf{Anonymisation} : Les données envoyées aux modèles Cloud (Gemini) sont débarrassées des informations personnellement identifiables.
    \item \textbf{Droit à l'accès} : L'administrateur peut extraire ou supprimer définitivement toutes les traces d'un étudiant (notes, messages, logs) sur simple demande.
\end{itemize}

\section{Conclusion}
L'analyse des cas d'utilisation a permis de valider la faisabilité des user stories définies dans le backlog. Nous passons maintenant à la conception de l'architecture permettant de supporter ces interactions.
